\section*{Chapitre 16 - Périmètre}
\addcontentsline{toc}{section}{Chapitre 16 - Périmètre}


\subsection*{I. Unités de longueur}
\addcontentsline{toc}{subsection}{I. Unités de longueur}


\begin{center}
\begin{tikzpicture}

  
  % Deuxième ligne - valeurs de position
  \node[draw, Couleur1,fill=Couleur1!20, minimum width=2.3cm, minimum height=0.8cm] at (-2,1.75) {\small Kilomètre};
  \node[draw, Couleur1,fill=Couleur1!20, minimum width=2.3cm, minimum height=0.8cm] at (0.3,1.75) {\small Hectomètre};
  \node[draw, Couleur1,fill=Couleur1!20, minimum width=2.3cm, minimum height=0.8cm] at (2.6,1.75) {\small Décamètre};
  \node[draw, Couleur4,fill=Couleur4!20, minimum width=2.3cm, minimum height=0.8cm] at (4.9,1.75) {\small Mètre};
  \node[draw, Couleur6,fill=Couleur6!20, minimum width=2.3cm, minimum height=0.8cm] at (7.2,1.75) {\small Décimètre};
  \node[draw, Couleur6,fill=Couleur6!20, minimum width=2.3cm, minimum height=0.8cm] at (9.5,1.75) {\small Centimètre};
  \node[draw, Couleur6,fill=Couleur6!20, minimum width=2.3cm, minimum height=0.8cm] at (11.8,1.75) {\small Millimètre};

  
  % Troisième ligne - noms des positions
  \node[draw, Couleur1, minimum width=2.3cm, minimum height=0.8cm, align=center] at (-2,0.95) {km};
  \node[draw, Couleur1, minimum width=2.3cm, minimum height=0.8cm, align=center] at (0.3,0.95) {hm};
  \node[draw, Couleur1, minimum width=2.3cm, minimum height=0.8cm, align=center] at (2.6,0.95) {dam};
  \node[draw, Couleur4, minimum width=2.3cm, minimum height=0.8cm, align=center] at (4.9,0.95) {m};
  \node[draw,Couleur6, minimum width=2.3cm, minimum height=0.8cm, align=center] at (7.2,0.95) {dm};
  \node[draw, Couleur6, minimum width=2.3cm, minimum height=0.8cm, align=center] at (9.5,0.95) {cm};
  \node[draw, Couleur6, minimum width=2.3cm, minimum height=0.8cm, align=center] at (11.8,0.95) {mm};

  
  % Quatrième ligne - les chiffres du nombre
  \node[draw, Couleur1, minimum width=2.3cm, minimum height=0.8cm] at (-2,0.15) {};
  \node[draw, Couleur1, minimum width=2.3cm, minimum height=0.8cm] at (0.3,0.15) {};
  \node[draw, Couleur1, minimum width=2.3cm, minimum height=0.8cm] at (2.6,0.15) {};
  \node[draw, Couleur4, minimum width=2.3cm, minimum height=0.8cm] at (4.9,0.15) {};
  \node[draw, Couleur6, minimum width=2.3cm, minimum height=0.8cm] at (7.2,0.15) {};
  \node[draw, Couleur6, minimum width=2.3cm, minimum height=0.8cm] at (9.5,0.15) {};
  \node[draw, Couleur6, minimum width=2.3cm, minimum height=0.8cm] at (11.8,0.15) {};

    % Quatrième ligne - les chiffres du nombre
  \node[draw, Couleur1, minimum width=2.3cm, minimum height=0.8cm] at (-2,-0.65) {};
  \node[draw, Couleur1, minimum width=2.3cm, minimum height=0.8cm] at (0.3,-0.65) {};
  \node[draw, Couleur1, minimum width=2.3cm, minimum height=0.8cm] at (2.6,-0.65) {};
  \node[draw, Couleur4, minimum width=2.3cm, minimum height=0.8cm] at (4.9,-0.65) {};
  \node[draw, Couleur6, minimum width=2.3cm, minimum height=0.8cm] at (7.2,-0.65) {};
  \node[draw, Couleur6, minimum width=2.3cm, minimum height=0.8cm] at (9.5,-0.65) {};
  \node[draw, Couleur6, minimum width=2.3cm, minimum height=0.8cm] at (11.8,-0.65) {};

\end{tikzpicture}
\end{center}
\vspace{-0.3cm}
\begin{Exemples}
\vspace{-0.4cm}
\begin{multicols}{2}
\begin{itemize}[label=\textcolor{Couleur8}{$\bullet$}]
\item $2,7 $ km $= 2,7 \times 1000$ m $= 2\,700 $ m
\item $560$ m $= 560 \div 1000 $ km $= 0,56 $ km
\end{itemize}
\end{multicols}
\vspace{-0.4cm}
\end{Exemples}


\subsection*{II. Périmètre}
\addcontentsline{toc}{subsection}{II. Périmètre}

\begin{Définition}[c]
Le \textbf{périmètre} d'un polygone est la somme des longueurs de ses côtés.
\end{Définition}

\begin{Exemple}
\vspace{-0.3cm}
\begin{minipage}{0.55\textwidth}
Le périmètre du triangle ABC est  \\
$\text{AB + BC + AC = 4 cm + 1,3 cm + 3,7 cm = 9 cm}$.
\end{minipage} \hfill
\begin{minipage}{0.4\textwidth}
\begin{center}
\begin{tikzpicture}[scale=.7]
    % Droite orange (Couleur8) passant par B et A
    \draw[Couleur8, very thick] (-1.5,0) -- (3.5,0);
    \draw[Couleur8, very thick] (-1.5,0) -- (1.7,-2);
    \draw[Couleur8, very thick] (3.5,0) -- (1.7,-2);
    % Points B et A avec des croix
    \node[thick] at (1.7,-2) {$+$};
    \node[thick] at (-1.5,0) {$+$};
    \node[thick] at (3.5,0) {$+$};
    
    % Étiquettes des points
    \node at (1.9,-2.3) {C};
    \node at (-1.7,0.3) {A};
    \node at (3.7,0.3) {B};
    
    % Les étiquettes
    \node[black] at (1,0.3) {4 cm};
    \node[black] at (-1,-1.3) {3,7 cm};
    \node[black] at (3.8,-1.3) {1,3 cm};
\end{tikzpicture}
\end{center}
\end{minipage} 

\end{Exemple}


\begin{Propriété}[c] \textbf{Formulaire.}
\begin{itemize}[label=\textcolor{Couleur1}{$\bullet$}]
\item Le périmètre $P$ d'un rectangle de longueur $L$ et de largeur $l$ est donné par la formule :\\
 $P_{\text{Rectangle}} = 2 \times (L + l)$ ou $P_{\text{Rectangle}} = 2 \times L + 2 \times l$.

\item Le périmètre $P$ d'un carré de côté $c$ est donné par la formule :  $P_{\text{Carré}}=4 \times c$

\item Le périmètre $P$ d'un disque de diamètre $D$ est égal à la longueur du cercle délimitant ce disque. \\
On a donc $P_{\text{Disque}} = \pi \times D$  ou $P_{\text{Disque}} = 2 \times \pi \times r$ (où r est le rayon du disque).
\end{itemize}
\end{Propriété}

\begin{Exemple}
Le périmètre d'un disque de rayon 6 cm vaut : P = 2 × $\pi$ × 6 cm =12$\pi \simeq$ 37,7 cm.
\end{Exemple}


\begin{Exemple}
\begin{minipage}{0.6\textwidth}
On remarque que le contour de la figure est constitué de trois segments et d'un demi-cercle. 
\begin{align*}
P_{\text{Figure}}&= \textcolor{Couleur8}{7,7 + 5,8 + 7,7} + \textcolor{Couleur6}{P_{\text{Demi-disque}}}\\
&= \textcolor{Couleur8}{7,7 + 5,8 + 7,7} + \textcolor{Couleur6}{(2 \times 2,9 \times \pi) \div 2} \\
&= \textcolor{Couleur8}{21,2} + \textcolor{Couleur6}{2,9 \pi} \\
&\simeq 30,3 \text{ cm}
\end{align*}
\end{minipage} \hfill
\begin{minipage}{0.4\textwidth}
\begin{center}
\begin{tikzpicture}[baseline,baseline=(current bounding box.north),scale = 0.35]

    \tikzset{
      point/.style={
        thick,
        draw,
        cross out,
        inner sep=0pt,
        minimum width=5pt,
        minimum height=5pt,
      },
    }
    \clip (-1.2,-1.3) rectangle (15.219999999999999,8.7);
    	\draw[color={black},line width = 0.5] (10.62,8)--(10.22,8)--(10.22,7.6)--(10.62,7.6)--cycle;
	\draw[color={black},line width = 0.5] (10.62,0)--(10.22,0)--(10.22,0.4)--(10.62,0.4)--cycle;
	\draw[color={black},line width = 0.5] (0,8)--(0,7.6)--(0.4,7.6)--(0.4,8)--cycle;
	\draw[color={black},line width = 0.5] (0,0)--(0.4,0)--(0.4,0.4)--(0,0.4)--cycle;
    	\draw[color ={Couleur8},line width = 1] (0,0)--(0,8);
	\draw[color ={Couleur8},line width = 1] (0,8)--(10.62,8);
	\draw[color ={Couleur8},line width = 1] (0,0)--(10.62,0);
	\draw[color={Couleur6},line width = 1] (10.62,0) arc (-90:90:4) ;
	\draw[color ={black}, dashed ] (10.62,8)--(10.62,0);
	\draw[color ={black}, dashed ] (10.62,4)--(14.6,3.65);

	\draw [color={black}] (5.31,0) node[anchor = center,scale=0.5, rotate = 0] {/};
	\draw [color={black}] (5.31,8) node[anchor = center,scale=0.5, rotate = 0] {/};
	\draw[color ={black},line width = 0.625,opacity = 0.8] (10.43,8.19)--(10.81,7.81);\draw[color ={black},line width = 0.625,opacity = 0.8] (10.43,7.81)--(10.81,8.19);\draw[color ={black},line width = 0.625,opacity = 0.8] (10.43,0.19)--(10.81,-0.19);\draw[color ={black},line width = 0.625,opacity = 0.8] (10.43,-0.19)--(10.81,0.19);

	\draw [color={black}] (12.61,3.83) node[anchor = center,scale=0.5, rotate = -5] {//};
	\draw [color={black}] (10.62,2) node[anchor = center,scale=0.5, rotate = -270] {//};
	\draw [color={black}] (10.62,6) node[anchor = center,scale=0.5, rotate = -90] {//};
	\draw [color={black}] (-0.7,4) node[anchor = center,scale=1.1, rotate = -90] {5,8 cm};
	\draw [color={black}] (5.31,0.7) node[anchor = center,scale=1.1, rotate = 0] {7,7 cm};
	\draw [color={black}] (12.67,4.53) node[anchor = center,scale=1.1, rotate = -5.02565] {{\small 2,9 cm}};

\end{tikzpicture}
\end{center}
\end{minipage} 

\end{Exemple}

